\documentclass[11pt,oneside]{article}

\pdfminorversion=7
\usepackage{QuizStyle}

\hypersetup{
    pdftitle={20261126 Quiz for MATH102 Calculus II},
    pdfauthor={Jungwoo Kim},
    pdfsubject={Quiz for MATH102 Calculus II},
    pdfcreator={LaTeX}
}

\begin{document}

\quiztitle{20261126 Quiz}

\begin{notation}
We use spherical coordinates $(\rho, \theta, \phi)$, where $\rho\ge0$,
$\theta\in[0, \pi]$ is the polar angle measured from the positive $z$-axis,
and $\phi\in[0, 2\pi]$ is the azimuthal angle measured counterclockwise
from the positive $x$-axis in the $xy$-plane. The textbook interchanges
$\theta$ and $\phi$.
\end{notation}

\begin{problem}{6.58}
Find the average value of the function $f(x, y, z)=x^2+y^2-z^2$
on the sets
\begin{problemsubparts}
    \item $[-1, 1]^3$.
    \item $x^2+y^2+z^2\le1$.
\end{problemsubparts}
\end{problem}
\begin{solution}
For a region $D$ of positive volume, the average value is
\[
    f_{\mathrm{avg}}=\frac{1}{\operatorname{Vol}(D)}
    \iiint_D f(x, y, z)\,\dd V.
\]
Both regions are invariant under permutations of the coordinates, so
by symmetry
\[
    \iiint_D x^2\,\dd V
    =\iiint_D y^2\,\dd V
    =\iiint_D z^2\,\dd V.
\]
Thus, by linearity,
\[
    \iiint_D f\,\dd V
    =\iiint_D(x^2+y^2-z^2)\,\dd V
    =\iiint_D x^2\,\dd V.
\]
\begin{problemsubparts}
    \item For $D=[-1, 1]^3$, we have $\operatorname{Vol}(D)=8$. Hence
    \[
    \begin{aligned}
        f_{\mathrm{avg}}
        &=\frac18\left(\int_{-1}^1x^2\,\dd x\right)
          \left(\int_{-1}^1\dd y\right)
          \left(\int_{-1}^1\dd z\right)\\
        &=\frac18\cdot\frac23\cdot2\cdot2
         =\boxed{\frac13}.
    \end{aligned}
    \]
    \item Let $D$ be the unit ball. By symmetry,
    \[
        \iiint_D x^2\,\dd V
        =\frac13\iiint_D(x^2+y^2+z^2)\,\dd V.
    \]
    In spherical coordinates,
    \[
        x^2+y^2+z^2=\rho^2, \qquad
        \dd V=\rho^2\sin\theta\,\dd\rho\,\dd\theta\,\dd\phi,
    \]
    where $0\le\rho\le1$, $0\le\theta\le\pi$, and $0\le\phi\le2\pi$.
    Since $\operatorname{Vol}(D)=4\pi/3$,
    \[
    \begin{aligned}
        f_{\mathrm{avg}}
        &=\frac{1}{3(4\pi/3)}
          \int_0^{2\pi}\int_0^\pi\int_0^1
          \rho^4\sin\theta\,\dd\rho\,\dd\theta\,\dd\phi\\
        &=\frac{1}{4\pi}\cdot\frac15\cdot2\cdot2\pi
         =\boxed{\frac15}.
    \end{aligned}
    \]
\end{problemsubparts}
\end{solution}

\begin{problem}{6.59}
Evaluate
\[
    \iiint_D xz^2\,\dd V
\]
for each of the following regions.
\begin{problemsubparts}
    \item The rectangular region $D=[1, 2]\times[3, 5]\times[-1, 10]$.
    \item The subset of $D=[1, 2]\times[3, 5]\times[-1, 10]$
          where $z>x+y$.
\end{problemsubparts}
\end{problem}
\begin{solution}
\begin{problemsubparts}
    \item Since $D$ is a rectangular box, the integral separates:
    \[
    \begin{aligned}
        \iiint_D xz^2\,\dd V
        &=\left(\int_1^2x\,\dd x\right)
          \left(\int_3^5\dd y\right)
          \left(\int_{-1}^{10}z^2\,\dd z\right)\\
        &=\frac32\cdot2\cdot\frac{1001}{3}
         =\boxed{1001}.
    \end{aligned}
    \]
    \item Let $D'=\{(x, y, z)\in D\mid z>x+y\}$. Since
    $4\le x+y\le7$ for $(x, y)\in[1, 2]\times[3, 5]$,
    the $z$-bounds are $x+y<z\le10$. The boundary $z=x+y$ has
    zero volume and does not affect the integral. Integrating successively
    and expanding the polynomial, we obtain
    \[
    \begin{aligned}
        \iiint_{D'} xz^2\,\dd V
        &=\int_1^2\int_3^5\int_{x+y}^{10}
          xz^2\,\dd z\,\dd y\,\dd x\\
        &=\frac13\int_1^2\int_3^5
          x\bigl(1000-(x+y)^3\bigr)\,\dd y\,\dd x\\
        &=\frac13\int_1^2x\left(
          2000-\frac{(x+5)^4-(x+3)^4}{4}\right)\dd x\\
        &=\int_1^2\left(-\frac23x^4-8x^3
          -\frac{98}{3}x^2+\frac{1864}{3}x\right)\dd x\\
        &=\left[-\frac{2}{15}x^5-2x^4
          -\frac{98}{9}x^3+\frac{932}{3}x^2\right]_1^2\\
        &=\boxed{\frac{36974}{45}}.
    \end{aligned}
    \]
\end{problemsubparts}
\end{solution}

\begin{problem}{6.60}
Let $D$ be the cube where $x$, $y$, and $z$ each vary from $-2$ to $2$.
\begin{problemsubparts}
    \item Explain by a symmetry argument why
          $\displaystyle\iiint_D xz^2\,\dd V=0$.
    \item Evaluate
          $\displaystyle\iiint_D(8x+2xz^2-4y^2z+10)\,\dd V$.
\end{problemsubparts}
\end{problem}
\begin{solution}
\begin{problemsubparts}
    \item For fixed $y$ and $z$, the function $xz^2$ is odd in $x$.
    Since $D=[-2, 2]^3$ is symmetric under $x\mapsto-x$,
    \[
        \iiint_D xz^2\,\dd V=\boxed{0}.
    \]
    \item The terms $8x$ and $2xz^2$ are odd in $x$, and
    $-4y^2z$ is odd in $z$. Their integrals vanish by symmetry.
    Since $\operatorname{Vol}(D)=4^3=64$,
    \[
    \begin{aligned}
        \iiint_D(8x+2xz^2-4y^2z+10)\,\dd V
        &=\iiint_D10\,\dd V\\
        &=10\cdot64=\boxed{640}.
    \end{aligned}
    \]
\end{problemsubparts}
\end{solution}

\begin{problem}{6.61}
A planet is represented by a ball of radius $R$ centered at the origin.
Its energy density [energy/volume] at $(x, y, z)$ is
\[
    f(x, y, z)=e^{-\sqrt{x^2+y^2+z^2}}.
\]
Find the total energy of the planet.
\end{problem}
\begin{solution}
Let $D=\{(x, y, z)\in\R^3\mid x^2+y^2+z^2\le R^2\}$.
The total energy is the integral of the energy density over $D$.
In spherical coordinates, $\sqrt{x^2+y^2+z^2}=\rho$, with
$0\le\rho\le R$, $0\le\theta\le\pi$, and $0\le\phi\le2\pi$.
Thus
\[
\begin{aligned}
    E
    &=\iiint_D f(x, y, z)\,\dd V\\
    &=\int_0^{2\pi}\int_0^\pi\int_0^R
      e^{-\rho}\rho^2\sin\theta\,\dd\rho\,\dd\theta\,\dd\phi\\
    &=\left(\int_0^{2\pi}\dd\phi\right)
      \left(\int_0^\pi\sin\theta\,\dd\theta\right)
      \int_0^R\rho^2e^{-\rho}\,\dd\rho\\
    &=4\pi\int_0^R\rho^2e^{-\rho}\,\dd\rho.
\end{aligned}
\]
Integration by parts gives
\[
\begin{aligned}
    \int\rho^2e^{-\rho}\,\dd\rho
    &=-\rho^2e^{-\rho}+2\int\rho e^{-\rho}\,\dd\rho\\
    &=-\rho^2e^{-\rho}-2(\rho+1)e^{-\rho}+C\\
    &=-(\rho^2+2\rho+2)e^{-\rho}+C.
\end{aligned}
\]
Therefore
\[
\begin{aligned}
    E
    &=4\pi\left[-(\rho^2+2\rho+2)e^{-\rho}\right]_0^R\\
    &=\boxed{4\pi\left(2-(R^2+2R+2)e^{-R}\right)}.
\end{aligned}
\]
\end{solution}

\begin{problem}{6.62}
Verify that the Jacobian determinant of the spherical-coordinate mapping
\[
    \mathbf F(\rho, \theta, \phi)=
    \begin{pmatrix}
        \rho\sin\theta\cos\phi\\
        \rho\sin\theta\sin\phi\\
        \rho\cos\theta
    \end{pmatrix}
\]
is $\det D\mathbf F(\rho, \theta, \phi)=\rho^2\sin\theta$.
\end{problem}
\begin{solution}
Taking partial derivatives in the order $(\rho, \theta, \phi)$,
we obtain
\[
    D\mathbf F(\rho, \theta, \phi)=
    \begin{pmatrix}
        \sin\theta\cos\phi & \rho\cos\theta\cos\phi
            & -\rho\sin\theta\sin\phi\\
        \sin\theta\sin\phi & \rho\cos\theta\sin\phi
            & \rho\sin\theta\cos\phi\\
        \cos\theta & -\rho\sin\theta & 0
    \end{pmatrix}.
\]
Factoring $\rho$ from the second column and $\rho\sin\theta$ from
the third column, we obtain
\[
    \det D\mathbf F(\rho, \theta, \phi)
    =\rho^2\sin\theta
    \begin{vmatrix}
        \sin\theta\cos\phi & \cos\theta\cos\phi & -\sin\phi\\
        \sin\theta\sin\phi & \cos\theta\sin\phi & \cos\phi\\
        \cos\theta & -\sin\theta & 0
    \end{vmatrix}.
\]
Expanding the remaining determinant along its third column gives
\[
\begin{aligned}
    &-\sin\phi
      \begin{vmatrix}
        \sin\theta\sin\phi & \cos\theta\sin\phi\\
        \cos\theta & -\sin\theta
      \end{vmatrix}
      -\cos\phi
      \begin{vmatrix}
        \sin\theta\cos\phi & \cos\theta\cos\phi\\
        \cos\theta & -\sin\theta
      \end{vmatrix}\\
    &\qquad=\sin^2\phi(\sin^2\theta+\cos^2\theta)
      +\cos^2\phi(\sin^2\theta+\cos^2\theta)\\
    &\qquad=(\sin^2\theta+\cos^2\theta)(\sin^2\phi+\cos^2\phi)=1.
\end{aligned}
\]
Hence
\[
    \det D\mathbf F(\rho, \theta, \phi)
    =\boxed{\rho^2\sin\theta}.
\]
\end{solution}

\end{document}
